
\documentclass[10pt, letterpaper]{article}
\usepackage{../../../../template/template}


\setDocTitle{Verbale esterno 18/02/2026}
\setDocVersion{1.0.0}
\setDocData{18/02/2026}
\setDocIndex{ve\_2026\_02\_18}
\setDocState{2}
\setDocRecipient{2}

\begin{document}

\makeTitlePage

\newpage

\tableofcontents

\newpage

\makeMeetingInfoTable{18/02/2026}{16:00}{17:00}{via \textit{Microsoft Teams$_G$}}

% ------------------------
% Partecipanti
% ------------------------
\addParticipant{Filippo}{Venzo}{Responsabile}{P}
\addParticipant{Valeria}{Baleanu}{Progettista}{P}
\addParticipant{Luca}{Granziero}{Verificatore}{P}
\addParticipant{Christian}{Libralato}{Progettista}{P}
\addParticipant{Francesco}{Pasqual}{Verificatore}{P}
\addParticipant{Leonardo}{Pellizzon}{Progettista}{P}
\addParticipant{Giuseppe}{De Fina}{Amministratore}{P}
\makeMeetingParticipantsTable

\section{Ordine del giorno}
\begin{itemize}
  \item Riscontro revisione RTB;
  \item valutazione tecnologie backend;
  \item strategie di caching;
  \item pianificazione attività future;
\end{itemize}

\section{Approfondimento}

\subsection{Riscontro revisione RTB}
Il gruppo ha riferito alla proponente l'esito della revisione RTB. 

\subsection{Valutazione tecnologie backend}
In seguito alle criticità sollevate dal professor Cardin circa l'utilizzo di \textit{Express$_{G}$} ritenuto poco strutturato per un'architettura complessa, il gruppo ha discusso insieme alla proponente il passaggio a \textit{NestJS$_{G}$}. Questa scelta garantirebbe una struttura più rigorosa ed una gestione delle dipendenze aiutata da un \textit{IoC Container$_{G}$} intergrato.

\subsection{Strategie di caching}
È stata discussa la proposta di utilizzare PostgreSQL con campi \textit{JSONB$_{G}$} e indici \textit{GIN$_{G}$} per la gestione del caching, in alternativa ad un \textit{database} non relazionale. La Proponente ha confermato la validità tecnica di questo approccio, raccomandando però particolare attenzione alla persistenza dei dati e alla gestione delle tabelle \textit{unlogged} per ottimizzare le performance senza compromettere la stabilità in caso di crash.

\subsection{Pianificazione delle attività future}
La proponente ha suggerito di sfruttare il tempo d'attesa della seconda revisione della \textit{milestone} RTB per predisporre la \textit{repository} del progetto. L'obiettivo è anticipare la configurazione dei \textit{configuration item$_{G}$} e della base \textit{frontend}.


% ------------------------
% Decisioni
% ------------------------

\addDecision{Conferma dell'approccio PostgreSQL (JSONB/GIN) per il caching dei dati}

\makeDecisionTable

% ------------------------
% TODO
% ------------------------

\addTodo{-}{Preparazione \textit{repository}}{Gruppo \teamName}{28/02/2026}
\makeTodoTable

\end{document}
